\section{Affine terminal forests at high girth}
\label{sec:random-terminal-forest-girth}

The larger-marker estimate also applies to a union of terminal paths.
This gives a quantitative bound for an entire affine terminal-join class.
All logarithms have base two.

\begin{lemma}[Translation by a terminal forest]
\label{lem:random-terminal-forest-translation}
Let $N$ be a normalized $k$-pole, where $k\ge2$ is even, and let $T$ be its set of ports.
Let $F_0$ be a $T$-join with the fewest edges, and put $P=V(F_0)$.
Then $F_0$ consists of $k/2$ vertex-disjoint paths.
For a uniform $T$-join $J$ and a uniform even subgraph $\xi$ of $N$,
\[
 \Pr(J\text{ is a forest})
 \le\Pr(\xi-P\text{ is a forest}).
\]
Moreover, if every terminal pair has distance at most $\ell$, then
\[
 |P|\le\frac k2(\ell+1),\qquad
 \operatorname{comp}(N[P])\le\frac k2.
\]
\end{lemma}

\begin{proof}
A cycle in $F_0$ could be deleted while preserving its boundary, contradicting minimality.
At a port the selected degree is one; at every other vertex it is zero or two.
Thus $F_0$ is a union of $k/2$ disjoint terminal paths.
In particular, $P$ contains all ports.

Pair the terminals arbitrarily and choose a shortest path for each pair.
Their symmetric difference is a $T$-join with at most $k\ell/2$ edges.
Minimality gives $|E(F_0)|\le k\ell/2$, and hence
$|P|=|E(F_0)|+k/2\le k(\ell+1)/2$.
The induced graph on $P$ may join some of these paths, so it has at most $k/2$ components.

The map $\xi\mapsto F_0\mathbin\triangle\xi$ is a bijection from the even space to the affine $T$-join space.
Outside $P$, its image has exactly the same edges as $\xi$.
If the image is a forest, its restriction outside $P$ is a forest.
This proves the probability inequality.
\end{proof}

\begin{theorem}[The affine forest bound at high girth]
\label{thm:random-terminal-forest-girth}
There are absolute constants $C_*,M,R$ such that the following holds.
Let $N$ be a normalized even-$k$ pole of rank $h\ge R$.
Put $x=\log h$ and $z=\log x$.
Assume
\[
 2\le k\le\frac{2x}{z^2},
 \qquad
 \operatorname{girth}(N)\ge
 \left\lceil2z+2\log z+2\log k+C_*\right\rceil.
\]
Then the proportion of forests in its affine all-port join space satisfies
\[
 \frac{\#\{T\text{-join forests in }N\}}{2^h}
 \le\frac12+M\frac{x^{11/2}}{\sqrt h}.
\]
No condition on bridges or cyclic connectivity is imposed.
\end{theorem}

\begin{proof}
Take $a=1/(1000A)$ from Theorem~\ref{thm:random-larger-marker}, and choose
$C_*=\lceil2\log(1000/a)\rceil$.
Let
\[
 g_*=\lceil2z+2\log z+2\log k+C_*\rceil.
\]
Uniformly over the allowed values of $k$, for sufficiently large $h$,
\[
 g_*\le8z,\qquad
 2^{-g_*/2}\le\frac{a}{1000kxz},\qquad
 |V(N)|=2h+k-2\le3h.
\]
Moore-ball packing bounds every terminal distance by a number $\ell$ with
\[
 \ell+1<8z\left(k+1+\frac{3ah}{1000kxz}\right).
\]
The preceding lemma therefore gives a marker satisfying
\[
 |P|\le4kz(k+1)+\frac{12ah}{1000x}
 \le a\frac h x
\]
for sufficiently large $h$.
The last comparison is uniform, because $k\le2x/z^2$ makes its first term polynomial in $x$, while $h=2^x$.
Also $\operatorname{comp}(N[P])\le k/2\le x/z^2$, and $P$ contains every vertex of degree below three.

Apply Theorem~\ref{thm:random-larger-marker} and the translation inequality.
Its constant $M$ and a sufficiently large absolute threshold work uniformly over $k$.
\end{proof}

For each fixed even $k$, the girth threshold is
$2\log\log h+2\log\log\log h+O_k(1)$.
In this range, a full affine terminal-forest space is impossible at large rank.
This does not establish a uniform non-saturation gap at fixed girth or a small-cut characterization of full four-terminal forest spaces.

\begin{theorem}[One marker for every terminal sector]
\label{thm:random-all-terminal-sectors-girth}
There are absolute constants $C_*,M,R$ with the following property.
Let $N$ be a normalized $k$-pole of rank $h\ge R$, put $x=\log h$, $z=\log x$, and assume
\[
 2\le k\le x^2,\qquad
 \operatorname{girth}(N)\ge
 \left\lceil2z+2\log z+2\log k+C_*\right\rceil.
\]
For every even subset $S$ of its ports, the number $F_S(N)$ of forests with odd vertex set $S$ satisfies
\[
 F_S(N)\le
 \left(\frac12+M\frac{x^{11/2}}{\sqrt h}\right)2^h.
\]
The estimate is uniform over all sectors, including when $k$ grows with $h$.
\end{theorem}

\begin{proof}
Take the same $a$ as in Theorem~\ref{thm:random-larger-marker} and
$C_*=\lceil2\log(1000/a)\rceil$.
The integer girth lower bound $g_*$ in the statement satisfies, uniformly for $k\le x^2$,
\[
 g_*\le8z,\qquad 2^{-g_*/2}\le\frac{a}{1000kxz}
\]
once $h$ is sufficiently large.
Also $|V(N)|\le3h$.
Consequently every terminal distance is at most a number $\ell$ satisfying
\[
 \ell+1<8z(k+1)+\frac{24ah}{1000kx}.
\]

Fix one port and join it to every other port by a shortest path.
Their union is connected and has at most $(k-1)\ell+1$ vertices.
Let $Q$ be a spanning tree of this union and put $P=V(Q)$.
Then $P$ contains all ports and is connected in $N$, while
\[
 |P|\le k(\ell+1)
 <8kz(k+1)+\frac{24ah}{1000x}
 \le a\frac h x.
\]
The final comparison is uniform for $k\le x^2$.

For any even port set $S$, the tree $Q$ has an $S$-join $J_S$.
Translation by $J_S$ maps the even space bijectively to the affine $S$-join space.
Outside $P$, the translated and original states agree.
Thus an $S$-join forest gives an even state that is a forest outside $P$.
Apply Theorem~\ref{thm:random-larger-marker}.
The same marker and the same probability bound work for every $S$.
\end{proof}

The same proof gives a more general sufficient condition, without the restriction $k\le(\log h)^2$.
If an integer girth lower bound $g_0\ge3$ satisfies
\[
 k g_0\left(k+1+(2h+k-2)2^{-g_0/2}\right)
 \le a\frac{h}{\log h},
\]
then the conclusion of Theorem~\ref{thm:random-all-terminal-sectors-girth} holds.
Indeed, Moore-ball packing bounds $k(\ell+1)$ by the left side, and the same connected geodesic marker applies.

\begin{lemma}[The growing-port aggregate]
\label{lem:random-growing-port-aggregate}
Let $Y$ be a normalized $k$-pole of rank $j$, where $k\ge2$, and set $\rho=j+k-1$.
If $\rho$ is sufficiently large and
\[
 2k-2\le a\frac{\rho}{\log\rho},
\]
where $a$ is the constant in Theorem~\ref{thm:random-larger-marker}, then
\[
 \sum_{S\text{ even}}F_S(Y)
 \le
 \left(\frac12+M\frac{(\log\rho)^{11/2}}{\sqrt\rho}\right)2^\rho.
\]
\end{lemma}

\begin{proof}
Take a disjoint tree with $k$ leaves and degree three at every nonleaf vertex; for $k=2$, take a single edge.
Add one edge from each leaf to a distinct port of $Y$.
Mark the new tree vertices.
The resulting subcubic graph has rank $\rho$.
The tree has $2k-2$ vertices, is connected, and contains every vertex of degree below three.
Restriction of uniform even states to $Y$ covers each even terminal-join class uniformly, with its $2^j$ states.
The event of a forest outside the marked tree therefore has probability
$2^{-\rho}\sum_S F_S(Y)$.
Apply the larger-marker theorem.
\end{proof}

\begin{corollary}[A quarter-space bound with a sufficiently large smaller shore]
\label{cor:random-high-girth-large-shores}
There are absolute constants $C_1,C_2$ with the following property.
Let connected loopless cubic graphs have ranks $r\to\infty$ and girths
\[
 g\ge
 \left\lceil4\log\log r+2\log\log\log r+C_1\right\rceil.
\]
Suppose a minimum cyclic cut has order $k$ and smaller shore rank
\[
 j\ge C_2 k\log(k+2).
\]
Then, uniformly over these graphs,
\[
 c(G)\le\left(\frac14+o(1)\right)2^r.
\]
In particular, these graphs are strictly below the half bound at sufficiently large rank.
The same conclusion holds with the smaller-shore hypothesis replaced by
$j\to\infty$ and $2k-2\le a(j+k-1)/\log(j+k-1)$.
\end{corollary}

\begin{proof}
Write $h\ge j$ for the shore ranks.
Minimality makes both shores connected.
It also makes their cut endpoints distinct: a vertex with two cut incidences is a leaf in its cyclic shore, and moving it to the other shore reduces the cut without destroying either cycle.
Thus the shores are normalized $k$-poles, with
\[
 r=h+j+k-1.
\]

At these ranks, a shortest cycle has cyclic complement.
Indeed, since $g\ge5$, it is induced and no outside vertex meets it twice.
Every component outside it consequently has minimum degree at least two.
Its boundary is a cyclic cut of size $g$.
The cubic Moore bound therefore gives
\[
 k\le g\le3\log r
\]
for large $r$.

Apply Theorem~\ref{thm:random-all-terminal-sectors-girth} to the shore of rank $h$.
Since $h\ge(r-k+1)/2$ and $k\le3\log r$, one has $h\ge r/3$ and $k\le(\log h)^2$ eventually.
Since $h\le r$ and $k\le3\log r$, the girth threshold of that theorem is at most
\[
 4\log\log r+2\log\log\log r+C_*+2\log3+1.
\]
Choose $C_1$ above this constant.
It follows that every even terminal sector in that shore has at most
$(1/2+o(1))2^h$ forests, uniformly.

Choose the absolute constant $C_2$ large enough that
$j\ge C_2 k\log(k+2)$ implies
$2k-2\le a(j+k-1)/\log(j+k-1)$ for all $k\ge2$.
For example, take $C_2\ge16$ with $C_2/(\log(2C_2)+4)\ge1/a$.
The function $s/\log s$ is increasing in the relevant range.
At $\rho_0=C_2 k\log(k+2)+k-1$ one has
\[
 \frac{\rho_0}{\log\rho_0}
 \ge
 \frac{C_2 k\log(k+2)}{\log(2C_2)+2\log(k+2)}
 \ge\frac{2C_2 k}{\log(2C_2)+4}
 \ge\frac{2k}{a},
\]
which proves the assertion.
Moreover $j\to\infty$ uniformly: if $k$ tends to infinity this follows from the rank hypothesis, while if $k$ stays bounded the shore order $2j+k-2$ must be at least the girth, which tends to infinity.
The growing-port aggregate consequently bounds the other shore's total forest mass by
$(1/2+o(1))2^{j+k-1}$.
For each fixed forest in that shore, at most the entire matching terminal sector in the first shore can combine with it to form one cycle.
Every cycle crossing the cut has a unique such pair of forest restrictions.
Hence the number of crossing cycles is at most
\[
 (1/2+o(1))2^h
 (1/2+o(1))2^{j+k-1}
 = (1/4+o(1))2^r.
\]
Cycles internal to the two shores number at most $2^h+2^j$.
Their ratio to $2^r$ tends to zero uniformly because $h\ge r/3$ and $j\to\infty$.
This proves the claim.
\end{proof}

\begin{proposition}[An exact criterion for full terminal-forest spaces]
\label{prop:random-full-terminal-forest-criterion}
Let $N$ be a normalized even-$k$ pole and $T$ its set of ports.
Every $T$-join is a forest if and only if every cycle $C$ of $N$ disjoint from $T$ leaves a component of $N-V(C)$ containing an odd number of terminals.
\end{proposition}

\begin{proof}
If every component of $N-V(C)$ has an even number of terminals, choose a terminal join in each component.
Their union with $C$ is a cyclic $T$-join.

Conversely, in a cyclic $T$-join choose one of its cycles $C$.
It avoids every port, whose selected degree is one.
Every vertex on $C$ has selected degree two, since the available degree is three and its parity is even.
No selected edge leaves $C$.
The remaining selected edges therefore form terminal joins in the components of $N-V(C)$.
Each such component has an even number of terminals.
\end{proof}

The standard-library checker \texttt{check\_random\_terminal\_forest\_girth.py}
verifies the translation, the common geodesic marker, and the saturation criterion
on 247 normalized poles.
It checks 8,312 affine states and 3,206 even terminal sectors.
These finite identities corroborate the implementation; the estimates above are proved independently of the checks.
